\section{Banco de questões — Computação em Nuvem}

\subsection*{N1 — Fundamento competitivo}

\begin{questao}{\nivel{N1} 16.01 — Cebraspe/Certo ou Errado}
Em SaaS, o provedor gerencia a aplicação e a infraestrutura subjacente, mas o cliente continua responsável, entre outros pontos, por usuários, dados e configurações sob seu controle.
\end{questao}

\begin{questao}{\nivel{N1} 16.02 — Cebraspe/Certo ou Errado}
Elasticidade e escalabilidade são sinônimos perfeitos: ambas significam apenas aumentar verticalmente a capacidade de uma única máquina.
\end{questao}

\begin{questao}{\nivel{N1} 16.03 — Cebraspe/Certo ou Errado}
Serverless não significa ausência de servidores físicos; significa abstração da administração direta da infraestrutura subjacente pelo consumidor do serviço.
\end{questao}

\begin{questao}{\nivel{N1} 16.04 — Cebraspe/Certo ou Errado}
Uma conexão dedicada entre datacenter e nuvem elimina, por definição, a necessidade de criptografia de dados em trânsito.
\end{questao}

\begin{questao}{\nivel{N1} 16.05 — Cebraspe/Certo ou Errado}
Security groups e listas de controle de rede podem operar em escopos diferentes, de modo que detalhes de statefulness e avaliação de regras dependem do provedor e do serviço específico.
\end{questao}

\begin{questao}{\nivel{N1} 16.06 — Cebraspe/Certo ou Errado}
A adoção de IaC melhora reprodutibilidade e rastreabilidade, mas pode propagar em escala uma configuração incorreta se revisão e validação forem fracas.
\end{questao}

\begin{questao}{\nivel{N1} 16.07 — Cebraspe/Certo ou Errado}
Multicloud implica, automaticamente, maior disponibilidade, independentemente de dependências comuns de identidade, DNS, rede ou pipeline.
\end{questao}

\begin{questao}{\nivel{N1} 16.08 — Cebraspe/Certo ou Errado}
Durabilidade de objetos e disponibilidade do serviço de armazenamento são propriedades distintas.
\end{questao}

\begin{questao}{\nivel{N1} 16.09 — Cebraspe/Certo ou Errado}
FinOps busca apenas reduzir gastos de nuvem ao menor valor possível, ainda que isso prejudique desempenho e resiliência.
\end{questao}

\begin{questao}{\nivel{N1} 16.10 — Cebraspe/Certo ou Errado}
A LGPD não exige, de forma genérica, que todo dado pessoal permaneça fisicamente no Brasil; transferências internacionais dependem das hipóteses e mecanismos jurídicos aplicáveis.
\end{questao}

\subsection*{N2 — Aplicação}

\begin{questao}{\nivel{N2} 16.11 — FGV/modelo de serviço}
Uma equipe utiliza banco de dados gerenciado no qual o provedor administra sistema operacional, patching e mecanismo do banco, enquanto a equipe gerencia schemas, usuários e dados. O serviço é mais próximo de:
\begin{enumerate}[label=\Alph*)]
\item IaaS puro;
\item PaaS/DBaaS;
\item SaaS de usuário final;
\item colocation;
\item bare metal não gerenciado.
\end{enumerate}
\end{questao}

\begin{questao}{\nivel{N2} 16.12 — FCC/escalabilidade}
Uma aplicação dobra o número de instâncias atrás de balanceador durante picos e reduz após a demanda. O comportamento caracteriza principalmente:
\begin{enumerate}[label=\Alph*)]
\item escalabilidade vertical apenas;
\item elasticidade com escala horizontal;
\item alta durabilidade;
\item backup incremental;
\item consistência forte.
\end{enumerate}
\end{questao}

\begin{questao}{\nivel{N2} 16.13 — Cebraspe/Certo ou Errado}
Autoscaling baseado apenas em CPU pode ser inadequado para aplicação cujo gargalo real seja fila de mensagens ou latência de banco.
\end{questao}

\begin{questao}{\nivel{N2} 16.14 — FGV/serverless}
Uma função é disparada por eventos esporádicos, executa por segundos e permanece ociosa na maior parte do dia. A principal vantagem de serverless nesse cenário é:
\begin{enumerate}[label=\Alph*)]
\item eliminar qualquer custo;
\item alinhar execução/escala ao evento sem administrar servidores diretamente;
\item garantir estado local permanente;
\item remover limites de execução;
\item eliminar cold start.
\end{enumerate}
\end{questao}

\begin{questao}{\nivel{N2} 16.15 — Cebraspe/Certo ou Errado}
Cold start pode afetar latência de funções serverless, especialmente em cargas esparsas ou runtimes pesados.
\end{questao}

\begin{questao}{\nivel{N2} 16.16 — FCC/rede}
Uma instância em subnet privada precisa acessar repositórios na Internet para atualização, mas não deve receber conexões iniciadas da Internet. Um componente comum para esse desenho é:
\begin{enumerate}[label=\Alph*)]
\item NAT gateway/serviço equivalente para saída;
\item load balancer público obrigatório;
\item DNS autoritativo;
\item object storage;
\item WAF apenas.
\end{enumerate}
\end{questao}

\begin{questao}{\nivel{N2} 16.17 — FGV/IAM}
Uma conta de serviço possui chave estática válida por dois anos e permissão administrativa em toda a assinatura. A melhoria prioritária é:
\begin{enumerate}[label=\Alph*)]
\item manter a chave e ocultá-la em variável local;
\item usar identidade de workload/credencial temporária e restringir permissões;
\item aumentar a validade da chave;
\item compartilhar a chave entre aplicações;
\item remover logs de acesso.
\end{enumerate}
\end{questao}

\begin{questao}{\nivel{N2} 16.18 — Cebraspe/Certo ou Errado}
MFA em contas administrativas não substitui privilégio mínimo e revisão periódica de acessos.
\end{questao}

\begin{questao}{\nivel{N2} 16.19 — FCC/KMS}
Na envelope encryption, os dados são normalmente cifrados por:
\begin{enumerate}[label=\Alph*)]
\item chave de dados, que é protegida por chave mestra/KEK;
\item certificado sem chave;
\item hash irreversível;
\item senha do usuário final;
\item regra de firewall.
\end{enumerate}
\end{questao}

\begin{questao}{\nivel{N2} 16.20 — FGV/FinOps}
Quarenta VMs de laboratório custam R\$ 0,80/h e permanecem ligadas 24x7 por 30 dias. O custo aproximado é:
\begin{enumerate}[label=\Alph*)]
\item R\$ 2.304;
\item R\$ 9.600;
\item R\$ 23.040;
\item R\$ 57.600;
\item R\$ 230.400.
\end{enumerate}
\end{questao}

\begin{questao}{\nivel{N2} 16.21 — Cebraspe/Certo ou Errado}
Budgets e alertas podem informar desvios de custo sem necessariamente bloquear automaticamente o consumo que excedeu a meta.
\end{questao}

\begin{questao}{\nivel{N2} 16.22 — FCC/storage}
Para data lake de objetos grandes, logs e arquivos sem necessidade de montar sistema de arquivos tradicional, o tipo de storage mais natural é:
\begin{enumerate}[label=\Alph*)]
\item objeto;
\item bloco exclusivamente;
\item memória RAM;
\item SAN Fibre Channel obrigatoriamente;
\item cache local efêmero apenas.
\end{enumerate}
\end{questao}

\begin{questao}{\nivel{N2} 16.23 — FGV/migração}
Mover uma aplicação para máquinas virtuais em nuvem com poucas mudanças arquiteturais corresponde, tipicamente, a:
\begin{enumerate}[label=\Alph*)]
\item rehost;
\item refactor completo;
\item retire;
\item repurchase necessariamente;
\item retain.
\end{enumerate}
\end{questao}

\begin{questao}{\nivel{N2} 16.24 — Cebraspe/Certo ou Errado}
Rehost pode acelerar migração, mas não captura automaticamente benefícios de serviços nativos, elasticidade e modernização.
\end{questao}

\begin{questao}{\nivel{N2} 16.25 — FCC/observabilidade}
Em microsserviços distribuídos entre nuvem e datacenter, o mecanismo mais útil para seguir uma mesma requisição de ponta a ponta é:
\begin{enumerate}[label=\Alph*)]
\item trace/correlation ID;
\item apenas CPU média;
\item DNS TTL;
\item storage class;
\item tag de custo.
\end{enumerate}
\end{questao}

\subsection*{N3 — Integração de concurso}

\begin{questao}{\nivel{N3} 16.26 — FGV/assertivas}
Considere:

I. Região e zona de disponibilidade representam níveis diferentes de domínio de falha em muitos provedores.

II. Distribuir instâncias em zonas distintas pode reduzir risco de falha de datacenter.

III. Distribuir em múltiplas regiões elimina automaticamente dependências comuns de identidade e DNS.

Assinale a alternativa correta.
\begin{enumerate}[label=\Alph*)]
\item Apenas I.
\item Apenas II.
\item Apenas I e II.
\item Apenas II e III.
\item I, II e III.
\end{enumerate}
\end{questao}

\begin{questao}{\nivel{N3} 16.27 — Cebraspe/Certo ou Errado}
Uma arquitetura multirregional pode aumentar resiliência e, simultaneamente, elevar complexidade de consistência, replicação, custo e soberania de dados.
\end{questao}

\begin{questao}{\nivel{N3} 16.28 — FCC/conectividade híbrida}
Uma organização usa circuito dedicado entre datacenter e nuvem, mas precisa proteger confidencialidade fim a fim. A conclusão correta é:
\begin{enumerate}[label=\Alph*)]
\item circuito dedicado torna criptografia desnecessária;
\item conectividade dedicada e criptografia são propriedades distintas e podem ser combinadas;
\item TLS não funciona em link privado;
\item VPN só pode operar na Internet pública;
\item Direct Connect/ExpressRoute implicam criptografia de aplicação por definição.
\end{enumerate}
\end{questao}

\begin{questao}{\nivel{N3} 16.29 — FGV/Zero Trust}
Uma empresa permite acesso administrativo a qualquer host que esteja na subnet de gestão. O princípio mais diretamente violado é:
\begin{enumerate}[label=\Alph*)]
\item confiar explicitamente em identidade e contexto em vez de localização;
\item escalabilidade horizontal;
\item object storage;
\item autoscaling;
\item caching.
\end{enumerate}
\end{questao}

\begin{questao}{\nivel{N3} 16.30 — Cebraspe/Certo ou Errado}
IaC versionada sem revisão de código, policy-as-code ou proteção de estado pode tornar erros de configuração mais rápidos e reproduzíveis, porém não necessariamente mais seguros.
\end{questao}

\begin{questao}{\nivel{N3} 16.31 — FCC/Terraform}
Uma equipe altera recursos manualmente no console após o deploy por Terraform. O risco relevante é:
\begin{enumerate}[label=\Alph*)]
\item configuration drift entre estado desejado e ambiente real;
\item impossibilidade de usar APIs;
\item perda automática de criptografia;
\item aumento de durabilidade;
\item ausência de logging.
\end{enumerate}
\end{questao}

\begin{questao}{\nivel{N3} 16.32 — FGV/segredos}
Uma imagem de contêiner publicada em registry privado contém credencial de banco em uma camada antiga. Mesmo que a última camada apague o arquivo, o risco persiste porque:
\begin{enumerate}[label=\Alph*)]
\item camadas anteriores podem continuar armazenando o segredo;
\item registry privado elimina histórico;
\item Docker converte segredo em hash;
\item imagens não podem ser inspecionadas;
\item o problema existe apenas em registry público.
\end{enumerate}
\end{questao}

\begin{questao}{\nivel{N3} 16.33 — Cebraspe/Certo ou Errado}
Certificação de segurança do provedor não transfere automaticamente ao cliente a conformidade de suas identidades, dados, aplicações e configurações.
\end{questao}

\begin{questao}{\nivel{N3} 16.34 — FCC/custo}
Uma equipe reduz custo desligando réplicas redundantes de banco durante a noite. O consumo cai, mas o RTO deixa de ser cumprido em falha noturna. A avaliação FinOps correta é:
\begin{enumerate}[label=\Alph*)]
\item a otimização foi bem-sucedida porque o custo caiu;
\item FinOps deve otimizar valor, equilibrando custo, risco e requisitos de serviço;
\item RTO não interfere em custo;
\item deve-se remover o SLO;
\item alta disponibilidade nunca deve ser considerada em FinOps.
\end{enumerate}
\end{questao}

\begin{questao}{\nivel{N3} 16.35 — FGV/data lake}
Um data lake possui milhões de objetos, mas sem catálogo, owner, classificação ou regras de retenção. O maior risco é:
\begin{enumerate}[label=\Alph*)]
\item data swamp, acesso indevido e custos sem governança;
\item excesso de normalização;
\item baixa largura de banda apenas;
\item ausência de RAID;
\item impossibilidade de usar IA.
\end{enumerate}
\end{questao}

\begin{questao}{\nivel{N3} 16.36 — Cebraspe/Certo ou Errado}
Aumentar a janela de contexto de um serviço de IA gerenciado não resolve, por si só, requisitos de classificação, residência e proteção dos dados enviados ao provedor.
\end{questao}

\begin{questao}{\nivel{N3} 16.37 — FCC/migração}
Uma aplicação legada possui baixa utilização, alto custo de licença e substituto SaaS funcionalmente equivalente. Antes de decidir por rehost, deve-se avaliar:
\begin{enumerate}[label=\Alph*)]
\item repurchase/substituição e até retire, conforme necessidade e TCO;
\item apenas tamanho da VM;
\item somente latência;
\item obrigatoriamente containerização;
\item somente escolha da região.
\end{enumerate}
\end{questao}

\begin{questao}{\nivel{N3} 16.38 — FGV/multicloud}
Dois provedores são usados, mas ambos dependem do mesmo IdP corporativo. Uma falha do IdP impede acesso aos dois ambientes. O cenário demonstra:
\begin{enumerate}[label=\Alph*)]
\item causa comum não eliminada pelo multicloud;
\item independência completa;
\item erro de storage;
\item problema exclusivo de billing;
\item ausência de VPC.
\end{enumerate}
\end{questao}

\begin{questao}{\nivel{N3} 16.39 — Cebraspe/Certo ou Errado}
Tags de custo são úteis para alocação e accountability, mas só geram valor se houver convenção, obrigatoriedade e tratamento de recursos sem identificação.
\end{questao}

\begin{questao}{\nivel{N3} 16.40 — FCC/observabilidade}
Um serviço distribuído possui logs em cinco contas cloud diferentes, sem sincronização de tempo nem correlation ID. O maior problema operacional é:
\begin{enumerate}[label=\Alph*)]
\item dificuldade de reconstruir incidentes e seguir requisições;
\item excesso de criptografia;
\item impossibilidade de autoscaling;
\item perda de durabilidade;
\item aumento de SLA.
\end{enumerate}
\end{questao}

\subsection*{N4 — Auditoria / avançado}

\begin{questao}{\nivel{N4} 16.41 — Cebraspe/caso integrado}
Uma conta cloud possui 2.000 permissões diretas a usuários, chaves estáticas sem rotação, ausência de tags e recursos criados fora de IaC. Julgue o item:

A migração para outro provedor, sem corrigir IAM e governança, tende apenas a transportar os mesmos riscos para novo ambiente.
\end{questao}

\begin{questao}{\nivel{N4} 16.42 — FGV/cloud audit}
Um órgão usa SaaS crítico sem exportação testada, sem inventário de dados e sem cláusula clara de eliminação no término. O procedimento prioritário é:
\begin{enumerate}[label=\Alph*)]
\item aumentar a quantidade de usuários;
\item testar portabilidade/saída, mapear dados e verificar retenção/eliminação contratual;
\item revisar apenas o SLA de uptime;
\item confiar na certificação do provedor;
\item reduzir logs para economizar.
\end{enumerate}
\end{questao}

\begin{questao}{\nivel{N4} 16.43 — FCC/FinOps}
A fatura anual mostra 14\% de recursos sem owner, 9\% de ambientes de teste ligados 24x7 e 6\% de snapshots sem retenção. A recomendação mais adequada é:
\begin{enumerate}[label=\Alph*)]
\item aumentar o orçamento;
\item estabelecer ownership/tagging, desligamento programado, retenção e monitoramento de custo;
\item migrar automaticamente tudo para on-premises;
\item eliminar snapshots;
\item desativar alertas de billing.
\end{enumerate}
\end{questao}

\begin{questao}{\nivel{N4} 16.44 — Cebraspe/Certo ou Errado}
Uma arquitetura pode ser altamente disponível dentro de uma região e ainda falhar completamente diante de comprometimento da conta raiz ou do provedor de identidade, se esses componentes forem causas comuns.
\end{questao}

\begin{questao}{\nivel{N4} 16.45 — FGV/soberania}
Um órgão transfere dados pessoais para região estrangeira porque ``o provedor é certificado''. A avaliação correta é:
\begin{enumerate}[label=\Alph*)]
\item certificação elimina qualquer análise jurídica;
\item localização, transferência internacional, contrato, acesso e base jurídica precisam ser avaliados separadamente;
\item dados cifrados nunca estão sujeitos à LGPD;
\item soberania se resume ao país do datacenter;
\item usar multicloud resolve automaticamente a transferência.
\end{enumerate}
\end{questao}

\begin{questao}{\nivel{N4} 16.46 — Cebraspe/Certo ou Errado}
Um relatório SOC/ISO do provedor pode ser evidência relevante, mas o auditor ainda precisa avaliar controles sob responsabilidade do cliente no modelo compartilhado.
\end{questao}

\begin{questao}{\nivel{N4} 16.47 — FCC/IaC}
Um pipeline Terraform aplica mudanças automaticamente com credencial administrativa, sem revisão, sem plano salvo e sem proteção do state. O principal conjunto de riscos é:
\begin{enumerate}[label=\Alph*)]
\item alteração não autorizada, segredo excessivo e comprometimento do estado/configuração;
\item apenas baixa disponibilidade de DNS;
\item falta de storage;
\item ausência de serverless;
\item excesso de versionamento.
\end{enumerate}
\end{questao}

\begin{questao}{\nivel{N4} 16.48 — FGV/continuidade}
Sistema em nuvem possui RTO de 1 h. O banco é replicado entre zonas, mas backups ficam apenas na mesma conta e a conta pode ser excluída pelo mesmo administrador que opera produção. O achado deve focar:
\begin{enumerate}[label=\Alph*)]
\item inexistência de replicação;
\item risco de causa comum administrativa e necessidade de proteção/segregação de backups e contas;
\item apenas performance;
\item ausência de CDN;
\item necessidade de multicloud obrigatória.
\end{enumerate}
\end{questao}

\begin{questao}{\nivel{N4} 16.49 — Cebraspe/Certo ou Errado}
A elasticidade pode aumentar desperdício se recursos crescerem automaticamente e não houver limites, desligamento, quotas ou governança de custo.
\end{questao}

\begin{questao}{\nivel{N4} 16.50 — FGV/matriz de achado}
A auditoria encontra 300 buckets, 41 públicos, 27 sem criptografia configurada e 80 sem owner. Qual achado é mais bem formulado?
\begin{enumerate}[label=\Alph*)]
\item ``Object storage é inseguro'';
\item ``A ausência de baseline, ownership e controle de configuração expõe dados e dificulta accountability; recomenda-se políticas preventivas, inventário, classificação, criptografia e revisão de exposição'';
\item ``Todos os buckets devem ser apagados'';
\item ``Basta mudar de provedor'';
\item ``O problema é somente custo''.
\end{enumerate}
\end{questao}
